
\chapter{Plan für die MA: Inhaltsverzeichnis}

\begin{enumerate}
	\item 
	Confidence intervals einführen
	\begin{enumerate}
		\item Definitionen
		\item asymptotic vs exact
		\item Bootstrapping
		\item insgesamt nur kurzes kapitel, da später die meisten methoden eingeführt werden wo direkt zensierte Daten mit ins Spiel kommen
	\end{enumerate}
	\item 
	Censoring einführen
	\begin{enumerate}
		\item Type I, Type II, hybrides Type I, hybrides Type II
		\item Probleme mit zensierten Daten erklären
	\end{enumerate}
	\item 
	CI for censored data, verschiedene Methoden
	\begin{enumerate}
		\item Pivoting the cdf
		\item Fiducial distribution
		\item Data Generating function
		\item Sundberg (Approximate, Binomial Inversionx)
	\end{enumerate}
\end{enumerate}

Wie kann ich DGF auf Hybrid zensierte Daten anwenden? Mehrere $c_i$ speichern? Das implementieren


\begin{table}[htbp]
	\centering
	\begin{tabular}{lll}
		\toprule
		Author & Censoring & Method/Model \\
		\midrule
		Jeng          & I               &                                  \\
		Sundberg      & I               & Exponential                      \\
		Wright        & I               & Two Parameter Exponential        \\
		Cramer2014a   & alle            & Pivoting CDF for Exponential     \\
		Lindqvist     & I, hybrid type II & Simulation Exponential        \\
		Cramer2024    & alle            &                                  \\
		\bottomrule
	\end{tabular}
	\caption{Paper im Überblick}
\end{table}




\chapter{TODO}

\begin{itemize}
	\item Binomial Inversion (Sundberg p. 396), erklären warum Binomial inversion nur gut bei Heavy censoring und Beispiel gebens
	\item Mehr Bootstrapping
	\item CI miteinander vergleichen (coverage prob, Laufzeit)
	\item C7 von Sundberg machen und verstehen
	\item 
	die Notation aus Cramer $\Upsilon$ benutzen, wo $\Upsilon(\mathbf D) = \max(D)$ z. b. 
	\item 
	Pivoting gut erklären. Lösung der Gleichungen 
	\[
	F^{\hat \vartheta}_{\vartheta}(\hat \vartheta_{\text{obs}}) = \alpha / 2
	\quad
	\text{ und }
	\quad
	F^{\hat \vartheta}_{\vartheta}(\hat \vartheta_{\text{obs}}) = 1 - \alpha / 2
	\]
	bzw. für one sided CI
	\[
	F^{\hat \vartheta}_{\vartheta}(\hat 	\vartheta_{\text{obs}}) = \alpha
	\quad
	\text{ oder }
	\quad
	F^{\hat \vartheta}_{\vartheta}(\hat \vartheta_{\text{obs}}) = 1 - \alpha
	\]
	\item 
	renv in R benutzen (Library Manager)
	\item 
	Notation TTT verwenden, nicht $S$
	\item 
	Kritik an Sundberg, z. B. dass er den Fisher Info Term für Type I einfach weglässt oder das Integral was mega dumm approximiert ist
	\item 
	die fisher info für uncensored type I and type II vergleichen und sagen, dass diese mit $c$ größer zunimmt und mit $m$ größer zunimmt, wie man an den Formeln sehen kann
	\item 
	Sundberg: Ich habe jetzt auch ein CI für type I censoring selber erstellt (das hat nicht die Form $[c_l \hat \theta, c_u \hat \theta]$). Vergleiche das mit dem von Sundberg.
	Für C2:
	\[
		1 - \alpha \approx \mathbb{P} \left( 
			z_{\alpha / 2} \leq \sqrt{n (1 - \exp(-c \theta))} \left( \frac{\hat \theta}{\theta} - 1 \right) \leq z_{1 - \alpha / 2}
		\right)
	\]
	\item 
	Laufzeiten messen (auch von neuen C1 und C2 exaktes für type I)
	\item 
	In der compare.py in Sundberg auch das dritte CI einbinden wo R immer eingegebn wird und in plots als grüne linie hinzufügen
	\item 
	die Tabelle für Beweis a, b in den Anhang
	\item 
	vor dem replizieren von Sundberg sagen, dass ich parametrisierung $1 / \theta$ wähle um konsistent mit Sundberg zu bleiben
\end{itemize}
